% !PRINTED 777
\weberpage{777}

\begin{center}
II.

\gesperrt{Die Irreducibilität der Kreistheilungsgleichung
und der Satz über die in einer Linearform enthaltenen
Primzahlen.}
\end{center}

Während im ersten Nachtrage, der als Ergänzung zum ersten Bande zu betrachten
ist, absichtlich kein Gebrauch von der Theorie der ganzen algebraischen Zahlen
gemacht ist, geben wir zum Schluss noch einen auf anderen Grundlagen, nämlich
auf den Sätzen über die Classenzahl (\S.\,192) beruhenden Beweis für die
Irreducibilität der Kreistheilungsgleichung, der wegen der Verallgemeinerungen,
die er zulässt, merkwürdig ist.

Es sei $m$ eine beliebige natürliche Zahl, und $n$ sei das System der zu $m$
theilerfremden positiven Zahlen; unter diesen giebt es
\weq{1}{\varphi(m) = \mu,}
die kleiner als $m$ sind, die wir mit $a$ bezeichnen.

Wenn wir alle nach dem Modul $m$ mit einem $a$ congruenten Zahlen $n$ in eine
Classe $A$ vereinigen, so erhalten wir $\mu$ Zahlclassen:
\weq{2}{A_1, A_2, \ldots, A_\mu,}
die bei der Composition durch Multiplication eine \gesperrt{Abel}'sche Gruppe
$\mu^{\text{ten}}$ Grades, $\mathfrak{N}$, bilden. Diese Gruppe ist schon im
\S.\,16 dieses Bandes betrachtet.

Die $\mu$ Charaktere dieser Gruppe bezeichnen wir mit
\weq{3}{\chi_1, \chi_2, \ldots, \chi_\mu,}
und setzen, wenn $\chi$ einer dieser Charaktere ist, und $n$ eine Zahl aus der
Classe $A$ bedeutet,
\weq{4}{\chi(A) = \chi(n).}
Diese Charaktere, die im \S.\,16 näher bestimmt sind, sind sämmtlich
$\mu^{\text{te}}$ Einheitswurzeln.
